\section{Introduction}
\label{sec:introduction}

Decentralised finance (DeFi) has rapidly evolved into a densely interconnected financial system in which protocols are linked through tokenised liabilities, rehypothecation, wrapping, bridging, and composable yield strategies. These linkages form an evolving \emph{credit-exposure network}: shocks to one protocol (e.g., insolvency, de-pegging, oracle failures, liquidity runs, or governance attacks) can propagate through downstream balance sheets and induce systemic stress. In such a setting, timely \emph{risk monitoring} and \emph{decision recommendation} are essential for investors, protocol operators, and regulators. Yet DeFi risk management remains challenging in practice: exposures are high-dimensional, time-varying, and often implicit; data are noisy and incomplete; and crises unfold quickly, leaving limited time for manual assessment.

A natural response is to build \emph{agentic} pipelines that automate surveillance and intervention planning. In principle, an agent can ingest on-chain measurements, maintain a belief over the evolving exposure network, surface early-warning signals, stress-test plausible shocks, and recommend risk-mitigating actions with auditable rationale. However, despite growing interest, agentic DeFi risk systems remain premature. A core bottleneck is the lack of a reliable \emph{foundation model backbone} for dynamic exposure networks: without accurate and data-efficient forecasting of graph time-series dynamics, monitoring agents are prone to false alarms and delayed detection, while decision agents risk recommending actions based on brittle or poorly calibrated beliefs. Unlike many recent multi-agent proposals that glue together large language models or other massive neural networks, our pipeline does not depend on such general-purpose big models; instead it uses a specialised graph–time‑series backbone tailored to the problem domain. In high-stakes contexts, uncertainty calibration and robustness to missingness and regime shifts are not optional - they determine whether an agent is trusted and deployable.

This paper addresses that bottleneck by building an agentic system around a strong graph time-series foundation model. We leverage DeXposure-FM, a recent foundation model trained on over $43.7$ million data entries, designed to forecast the evolution of DeFi exposure networks in a multi-horizon setting. On top of this predictive backbone, we develop \emph{DeXposure-Agent}, which, to our best knowledge, is the first agentic system for DeFi risk monitoring and decision recommendation. The system instantiates a conservative but operationally meaningful loop: it continuously ingests exposure snapshots, uses {DeXposure-FM} to forecast future credit-exposure graphs, computes systemic risk indicators and standardised stress scenarios, and emits auditable alerts and playbook-based recommendations under explicit data-health and uncertainty gates. The design is intentionally modular: the foundation model provides probabilistic forecasts; the monitoring layer translates forecasts into interpretable risk signals; the scenario engine evaluates stress outcomes; and the decision layer maps signals into recommendations while enforcing guardrails for safe deployment.

%\paragraph{DeXposure-Agent overview.}
DeXposure-Agent follows a \emph{forecast-measure-decision} paradigm. At each epoch, it (1) checks the integrity of the incoming graph snapshot via a data-health gate, (2) forecasts multi-step future exposure networks using {DeXposure-FM}, (3) computes risk measurements on predicted graphs (e.g., concentration, spillovers, and systemic importance) and evaluates a standardised library of stress scenarios, and (4) outputs a ranked set of alerts and decision tickets. Each alert is accompanied by evidence (triggering metric deviations), attribution (top contributing nodes/edges/sectors), and uncertainty-aware confidence. Each decision ticket is conservative and auditable: it links back to explicit triggers, includes scenario-based impact summaries, and is suppressed or downgraded when data quality or model reliability degrades. This design reflects a key principle: in DeFi risk management, the ability to say \emph{``we do not know''} is as important as the ability to recommend an action.

To make this pipeline concrete we implemented a modular architecture. on-chain TVL/token data are pulled from public APIs such as DefiLlama, transformed through a token–protocol mapping and assembled into discrete exposure graph snapshots; a lightweight data-health service scores each snapshot and triggers a safe mode when needed. The prediction backbone runs as a separate inference service—the DeXposure‑FM model returns probability distributions over future graphs along multiple horizons, supporting both expected‑weight construction and Monte Carlo sampling. A monitoring layer computes a set of systemic indices on the representative forecast, maintains rolling baselines for triggers, attributes alerts by marginal contribution, and quantifies forecast uncertainty. A parallel scenario engine applies a library of standardised shocks to the same forecasts and ranks outcomes by expected/tail loss. Finally a decision module consults a small playbook of conservative actions, enforces data‑health and confidence gates, scores and ranks tickets, and exposes results via an API and dashboard. All components communicate through simple data structures, so the system can be deployed as microservices or as a single Python application depending on operational needs.

%\paragraph{Evaluation and findings.}
We evaluate DeXposure-Agent across xx benchmarks and stress-test settings, including low-data and weakly supervised regimes that reflect real operational constraints. We compare against state-of-the-art methods for spatio-temporal/graph risk forecasting and against agentic baselines that do not leverage a dedicated graph time-series foundation model backbone. The results show that DeXposure-Agent: (1) produces earlier and more stable early-warning signals of systemic fragility,
(2) improves risk metric forecasting and uncertainty calibration, and
(3) yields higher-quality recommendation tickets (more robust, better justified, and less sensitive to noise),
all significantly outperforming existing approaches in both predictive accuracy and operational utility.
An open-source implementation is available at
\href{https://github.com/EVIEHub/DeXposure-Agent}{https://github.com/EVIEHub/DeXposure-Agent}.